\subsection{معماری و پیاده‌سازی}
پیاده‌سازی ماژول در فایل \lr{Codes/framework/compression\_engine/admm\_pruner.py} با ساختار صنعتی و مولفه‌های کلیدی زیر همراه است:
\begin{itemize}
\item \textbf{لایه خطی ماسک‌دار \lr{ADMMMaskedLinear}:} جایگزینی مستقیم لایه‌های خطی استاندارد با یک لایه مجهز به بافر ماسک غیرپایدار و اعمال تخمین‌گر مستقیم \lr{ADMMSTEPrune} جهت انتشار روان گرادیان‌ها در پس‌انتشار:
\begin{equation}
\mathbf{y} = \operatorname{Linear}\left(\mathbf{x}, \mathbf{W} \odot \mathbf{M}, \mathbf{b}\right)
\end{equation}
\item \textbf{بازکالیبراسیون نرمال‌سازی لایه \lr{(LayerNorm Recalibration)}:} تنظیم مجدد پارامترهای میانگین و واریانس در لایه‌های \lr{LayerNorm} مدل پس از اعمال ماسک، جهت اصلاح شیفت توزیع فعال‌سازی ناشی از حذف اتصالات.
\item \textbf{سازگاری کامل با سخت‌افزارهای تنسور کُر \lr{NVIDIA Ampere}:} استفاده از چیدمان ۲:۴ امکان بهره‌گیری حداکثری از شتاب دوبرابری پایپ‌لاین‌های محاسباتی ماتریسی مدرن را فراهم می‌سازد.
\end{itemize}

\subsection{محدودیت‌ها و تحلیل شکست}
مشاهدات آزمایشگاهی حاصل از ارزیابی هرس \lr{ADMM} ابعاد تحلیلی زیر را مشخص می‌سازد:
\begin{itemize}
\item \textbf{نرخ قبولی ۱۰۰ درصدی در کلیه مدل‌ها:} روش \lr{ADMMPruner} بر روی هر چهار معماری ناهمگن با موفقیت به وضعیت قبولی \lr{PASS} با سطح دقت بین \lr{78.21\%} تا \lr{84.29\%} نائل آمد.
\item \textbf{نقش تقطیر بازسازی لایه‌ای در حفظ دقت:} مقایسه عملکرد با روش‌های هرس مرسوم نشان می‌دهد که حضور جمله تقطیر $\mathcal{L}_{\text{distill}}$ و بازکالیبراسیون \lr{LayerNorm} نقش حیاتی در جلوگیری از فروپاشی شاخص F1 (بالای ۰٫۷۹ در تمام مدل‌ها) ایفا کرده است.
\item \textbf{تحلیل حساسیت به گام‌های کالیبراسیون:} در لایه‌های حساس مانند پیش‌بینی کلمات انتهایی، در صورتی که تعداد گام‌های تقطیر لایه‌ای (\lr{distill\_steps}) ناکافی باشد یا داده‌های ورودی کالیبراسیون تنوع دامنه زبانی لازم را پوشش ندهند، هرس ۲:۴ می‌تواند به اعوجاج الگوهای نحوی ظریف بیانجامد. همچنین مصرف حافظه در طول کالیبراسیون به دلیل فعال‌سازی ذخیره تانسورهای موقت نسبت به هرس ایستا اندکی بالاتر است.
\end{itemize}
